% authority: science-draft
% status: proposal-only
% route: ontology-law-research-packet
% trigger_classification: derivation_critical_missing_source_law
% adoption_status: blocked_adoption_open_continuation
\documentclass[11pt]{article}
\usepackage[margin=1in]{geometry}
\usepackage{amsmath,amssymb,amsthm}
\usepackage{booktabs}
\usepackage{enumitem}

\newtheorem{definition}{Definition}
\newtheorem{theorem}{Theorem}
\newtheorem{proposition}{Proposition}
\newtheorem{remark}{Remark}
\newcommand{\Ftwo}{\mathbb F_2}
\newcommand{\Disc}{\mathsf{Disc}_{\mathrm{src}}}
\newcommand{\AdmVar}{\mathsf{AdmVar}_{\mathrm{src}}}
\newcommand{\EqCand}{\mathsf{EqSrc}^{\mathrm{cand}}}
\newcommand{\LabelFix}{\mathsf{LabelFix}_{\chi}}
\newcommand{\RelPres}{\mathsf{RelPres}}
\newcommand{\RelRefl}{\mathsf{RelRefl}}
\newcommand{\RelAut}{\mathsf{RelAut}}
\newcommand{\Bottom}{\mathsf{Bottom}_{\mathrm{src}}}

\title{Proposal-Only EqSrc Intrinsic Discriminator and Admissibility Law Candidate v3:\\
Typed Cross-Complex Reflection-Scope Repair}
\author{Ontology Formalizer control packet RT-20260718-031}
\date{2026-07-19}

\begin{document}
\maketitle

\section{Control Status}

This packet defines
\(\mathsf{EqSrcIntrinsicDiscriminatorAdmissibilityLaw}^{\mathrm{cand},v3}_{\mathrm{src}}\)
as a proposal-only repair of the cross-complex typing and chain-map
reflection-scope defect identified by RT-20260718-030. Current ontology does
not derive this law, select finite chain packages, select their differentials,
select homology as a physical source discriminator, select a physical category
of chain maps, or establish boundary moves as physically admissible source
variations.

\begin{verbatim}
candidate_result: candidate_repaired_pending_fresh_smuggling_audit
candidate_name: EqSrcIntrinsicDiscriminatorAdmissibilityLaw_src^cand,v3
obstruction_addressed: OBST-EQSRC-INTRINSIC-DISCRIMINATOR-CHAIN-MAP-REFLECTION-001
obstruction_status: addressed_in_candidate_not_independently_cleared
formal_model: finite_2_truncated_chain_complex_over_F2
cross_complex_relation_preservation: explicitly_typed
cross_complex_relation_reflection: explicitly_typed
chain_map_preservation: always_under_displayed_hypotheses
chain_map_reflection: iff_induced_H1_is_injective
chain_isomorphism_status: sufficient_not_necessary
noninvertible_reflecting_witness_states: 4
noninvertible_reflecting_witness_pair_checks: 16
noninvertible_reflecting_witness_mismatches: 0
current_ontology_derives_candidate: false
physical_admissibility_established: false
physical_covariance_established: false
candidate_status: proposal-only
artifact_status: draft/control
source_extension_status: source-extension data
adoption_status: blocked_adoption_open_continuation
general_EqSrc_discharged: false
distance_to_gr_ledger_changed: false
metric_use_ledger_changed: false
freeze_decision: not_frozen
next_route: fresh_smuggling_auditor_review
\end{verbatim}

The exact v2 source remains historical evidence. This v3 packet repairs its
endomap-only notation and withdraws whole-chain noninvertibility as a
necessary failure condition for relation reflection. It preserves the
surviving v2 theorem core and does not adopt or promote the candidate.

\section{Supplied Source Packages}

\begin{definition}[Admitted chain package]
An admitted source package is a supplied finite two-truncated chain complex of
\(\Ftwo\)-vector spaces
\[
E=\left(C_2(E)\xrightarrow{\partial_2^E}
C_1(E)\xrightarrow{\partial_1^E}C_0(E)\right),
\qquad \partial_1^E\partial_2^E=0.
\]
Define
\[
A_E=Z_1(E)=\ker\partial_1^E,\qquad
B_E=B_1(E)=\operatorname{im}\partial_2^E,\qquad
K_E=H_1(E;\Ftwo)=A_E/B_E,
\]
\[
\Disc(E)=\chi_E:A_E\to K_E,\qquad \chi_E(z)=[z],
\]
\[
z\,\EqCand_E\,z'
\quad\Longleftrightarrow\quad
\chi_E(z)=\chi_E(z')
\quad\Longleftrightarrow\quad z-z'\in B_E.
\]
\end{definition}

The same definitions with primes apply to a second supplied source package
\(E'\). The prime distinguishes two objects in the candidate source category;
it does not denote a target manifold, empirical target, target atlas, target
metric, detector outcome, or downstream GR object.

Once \(E\) is supplied, these are conservative finite algebra. Current
ontology does not select the package family, field, truncation degree,
differentials, homology degree, discriminator, or physical meaning of the
relation.

\begin{definition}[Independently declared boundary moves]
Before using \(\chi_E\) or \(\EqCand_E\), declare
\[
T_b(z)=z+\partial_2^E b,\qquad b\in C_2(E),
\]
and let \(\AdmVar(E)\) be the transformation monoid generated by finite
compositions of the maps \(T_b\).
\end{definition}

Thus
\[
\AdmVar(E)=\{\tau_c:A_E\to A_E\mid \tau_c(z)=z+c,\ c\in B_E\}.
\]
This declaration does not establish that these moves are physically
admissible.

\section{Typed Cross-Complex Relation Predicates}

\begin{definition}[Typed preservation and reflection]
For supplied packages \(E,E'\) and a typed map \(G:A_E\to A_{E'}\), define
\[
\RelPres_{E,E'}(G)
\quad\Longleftrightarrow\quad
z\,\EqCand_E\,z'\Longrightarrow
Gz\,\EqCand_{E'}\,Gz'
\quad\text{for all }z,z'\in A_E,
\]
and
\[
\RelRefl_{E,E'}(G)
\quad\Longleftrightarrow\quad
Gz\,\EqCand_{E'}\,Gz'\Longrightarrow
z\,\EqCand_E\,z'
\quad\text{for all }z,z'\in A_E.
\]
\end{definition}

For \(E=E'\), these specialize to the self-map predicates used in v2. A
cross-complex pointwise label-fixing predicate is not defined here: the
equation \(\chi_{E'}G=\chi_E\) is ill typed without a separately supplied
comparison between \(K_E\) and \(K_{E'}\).

\begin{theorem}[Typed quotient factorization]
For any typed map \(G:A_E\to A_{E'}\),
\[
\RelPres_{E,E'}(G)
\quad\Longleftrightarrow\quad
\text{there exists a unique }\bar G:K_E\to K_{E'}
\text{ with }\chi_{E'}G=\bar G\chi_E.
\]
When this factorization exists,
\[
\RelRefl_{E,E'}(G)
\quad\Longleftrightarrow\quad
\bar G\text{ is injective}.
\]
\end{theorem}

\begin{proof}
If \(G\) preserves the relation, set
\(\bar G([z])=[Gz]\). Equal source classes are related, so their images are
related and the definition is well defined. Surjectivity of \(\chi_E\)
forces uniqueness. Conversely, the factorization maps equal source quotient
values to equal destination quotient values, hence preserves the relation.

Under the factorization,
\(\bar G([z])=\bar G([z'])\) is equivalent to
\(Gz\,\EqCand_{E'}\,Gz'\). Reflection is therefore exactly the implication
\([z]=[z']\), which is injectivity of \(\bar G\).
\end{proof}

\section{Chain-Map Reflection Theorem}

\begin{theorem}[Exact chain-map reflection criterion]
Let \(F:E\to E'\) be a chain map. Then \(F_1\) restricts to a typed map
\[
F_1|_{A_E}:A_E\longrightarrow A_{E'},
\]
which always preserves the candidate relation. It induces
\[
H_1(F):K_E\longrightarrow K_{E'}
\]
and obeys
\[
\chi_{E'}F_1|_{A_E}=H_1(F)\chi_E.
\]
Moreover,
\[
\RelRefl_{E,E'}(F_1|_{A_E})
\quad\Longleftrightarrow\quad
H_1(F)\text{ is injective},
\]
equivalently,
\[
(F_1|_{A_E})^{-1}(B_{E'})=B_E.
\]
\end{theorem}

\begin{proof}
The chain equation
\(\partial_1^{E'}F_1=F_0\partial_1^E\) maps cycles to cycles, and
\(F_1\partial_2^E=\partial_2^{E'}F_2\) maps boundaries into boundaries.
Consequently \([z]\mapsto[F_1z]\) is well defined and gives \(H_1(F)\) and
the displayed factorization. Preservation follows from
\[
z-z'\in B_E\Longrightarrow F_1(z-z')\in B_{E'}.
\]
Reflection is
\[
F_1(z-z')\in B_{E'}\Longrightarrow z-z'\in B_E,
\]
which is exactly \(\ker H_1(F)=0\), equivalently injectivity. Setting
\(z'=0\) gives the inverse-image formulation.
\end{proof}

\begin{proposition}[Chain isomorphism is sufficient, not necessary]
Every chain isomorphism induces a homology isomorphism and therefore preserves
and reflects the candidate relation. Whole-chain invertibility, degree-one
invertibility, and surjectivity of \(H_1(F)\) are not necessary for
reflection.
\end{proposition}

\begin{proof}
An inverse chain map induces the inverse homology map, so a chain isomorphism
gives an injective \(H_1(F)\). The noninvertible witness below has injective
induced homology and proves that the converse fails. An injective but
nonsurjective induced homology map also reflects the relation without
surjecting onto destination quotient classes.
\end{proof}

Relation reflection means only that quotient classes do not coalesce under
the supplied map. It is not physical reversibility, causal invertibility,
losslessness, observer independence, or physical covariance.

\section{Four-State Noninvertible Reflecting Witness}

Let
\[
C_2=\Ftwo\langle f\rangle,\qquad
C_1=\Ftwo\langle a,b\rangle,\qquad C_0=0,
\qquad \partial_2f=a.
\]
Then
\[
A=C_1,\qquad B=\{0,a\},\qquad K\cong\Ftwo\langle[b]\rangle,
\]
with relation classes
\[
\{0,a\},\qquad \{b,a+b\}.
\]

Define a self-chain-map by
\[
F_2=0,\qquad F_1(a)=0,\qquad F_1(b)=b,\qquad F_0=0.
\]
The chain condition holds because
\[
F_1\partial_2f=F_1(a)=0=\partial_2F_2f.
\]
Writing a state as \(\alpha a+\beta b\),
\[
F_1(\alpha a+\beta b)=\beta b.
\]
Two states are related exactly when their \(\beta\)-coordinates agree, and
their images are related under the same condition. Hence all 16 ordered-pair
tests agree: eight related pairs, eight unrelated pairs, and zero mismatches.

The image and kernel of \(F_1\) each have two elements, so \(F_1\) is
noninvertible. Nevertheless
\[
H_1(F)[b]=[b],
\]
so \(H_1(F)\) is the identity and is injective. This refutes whole-chain
invertibility as a necessary reflection condition without establishing any
physical interpretation of the chain map.

\begin{remark}[Degenerate edge cases]
If \(K_E=0\), every induced map is injective and reflection holds. If
\(K_{E'}=0\), reflection holds exactly when \(K_E=0\). These are quotient
facts, not physical completeness or triviality claims.
\end{remark}

\section{Surviving v2 Self-Map Core}

\begin{definition}[Self-map label and relation predicates]
For \(V:A_E\to A_E\), retain
\[
\LabelFix(V)\Longleftrightarrow
\chi_E(Vz)=\chi_E(z)\quad\text{for every }z\in A_E.
\]
Relation preservation and reflection are the \(E=E'\) cases above. A
relation automorphism is a bijective self-map with both predicates.
\end{definition}

Under the unique self-map factorization
\(\chi_EV=\bar V\chi_E\), pointwise label fixing is
\(\bar V=\operatorname{id}_{K_E}\), while relation reflection is only
injectivity of \(\bar V\). Labels are mathematical quotient values, not
established observables.

\begin{proposition}[Exact translation classification]
For \(\tau_c(z)=z+c\), \(c\in A_E\),
\[
\chi_E(\tau_cz)=\chi_E(z)+[c].
\]
Thus \(\tau_c\) fixes labels pointwise exactly when \(c\in B_E\), while every
translation is a relation automorphism because pairwise differences are
unchanged.
\end{proposition}

\begin{proposition}[Linear self-map criterion]
For a linear self-map \(L:A_E\to A_E\), relation preservation is equivalent
to \(L(B_E)\subseteq B_E\). For invertible \(L\), relation automorphism is
equivalent to \(L(B_E)=B_E\).
\end{proposition}

\begin{proposition}[Conditional quotient uniqueness]
If \(d:A_E\to D\) is surjective and
\[
d(z)=d(z')\quad\Longleftrightarrow\quad z-z'\in B_E,
\]
there is a unique bijection \(\bar d:K_E\to D\) with
\(d=\bar d\chi_E\).
\end{proposition}

\begin{proposition}[Finite admitted-move robustness]
Every finite word in \(\AdmVar(E)\) is translation by an element of \(B_E\).
It therefore fixes \(\chi_E\) pointwise and preserves and reflects the
candidate relation.
\end{proposition}

For the v2 eight-state package \(E_\star\), the corrected finite results
remain:
\[
A_\star=\Ftwo^3,\qquad B_\star=\{0,e_0+e_1\}.
\]
Translation by \(e_2\) changes all eight quotient labels and gives zero
relation mismatches across 64 ordered pairs. The ambient invertible linear map
\[
L(e_0)=e_0,\qquad L(e_1)=e_0+e_2,\qquad L(e_2)=e_1
\]
maps \(e_0+e_1\) outside \(B_\star\) and gives 16 relation mismatches. It is
not a boundary move, an admitted chain map, or an established physical
variation.

\section{Repaired Fail-Closed Branches}

The proposal returns \(\Bottom\), or withholds the named certificate, under
the following exact conditions:
\begin{enumerate}[leftmargin=*]
  \item return \(\Bottom\) if a supplied package fails
  \(\partial_1\partial_2=0\), or a displayed map is mistyped;
  \item reject a preservation claim if a related source pair maps to an
  unrelated destination pair, equivalently if the quotient factorization
  fails;
  \item reject a reflection claim if the induced quotient map is proved
  noninjective;
  \item for a chain map, reject reflection when \(H_1(F)\) is proved
  noninjective, equivalently when a cycle outside \(B_E\) maps into
  \(B_{E'}\);
  \item if injectivity has not been established, withhold the reflection
  certificate without asserting noninjectivity;
  \item do not reject reflection merely because \(F\), \(F_1\), or the whole
  chain package is noninvertible;
  \item do not assert a cross-complex pointwise label comparison without a
  separately typed codomain comparison;
  \item reject any overread of naturality as physical covariance, reflection
  as physical reversibility, a chain isomorphism as an adopted physical
  symmetry, or quotient labels as observables;
  \item reject physical selection inferred from the shared differential,
  definition order, a finite witness, a validator, a role, a registry,
  generated output, handoff, checkpoint, commit, or recency; and
  \item fail closed if target or workflow data enters the source-definition
  dependency closure.
\end{enumerate}

Failure closes only the stated candidate instance or certificate. It does not
establish a general EqSrc no-go, a program-wide no-go, or future
source-extension impossibility.

\section{Primitive Selection and Physical Scope}

The supplied differential remains the shared mathematical root
\[
\partial_2\longrightarrow B_E\longrightarrow\chi_E,\EqCand_E,
\qquad
\partial_2\longrightarrow\{T_b\}\longrightarrow\AdmVar(E).
\]
This alignment is transparent and noncircular as finite algebra, but it may
align relation and variation by construction. It does not establish
independent physical selection.

Current ontology does not derive or select the finite chain-package family,
field, truncation degree, differentials, homology degree, homology
discriminator, boundary translations, category of physically admissible
chain maps, physical covariance, continuum applicability, or robustness
beyond the declared finite hypotheses.

\section{No-Target and No-Process Guard}

The transitive mathematical dependency closure remains confined to displayed
finite source algebra. Desired exact-GR recovery is a later obligation, not a
premise for the chain package, discriminator, variation class, relation
predicates, induced homology, or witnesses.

No target atlas, target topology, target metric, detector outcome, benchmark
result, registry value, role identity, validator result, generated derivative,
handoff, checkpoint, commit, file order, or recency supplies a defining
scientific premise. Those surfaces establish provenance or transaction
integrity only.

\section{Distance-to-GR, Freeze, and Next Route}

This packet claims the exact cross-complex typing and reflection-scope defect
repaired at the proposal level, pending fresh independent audit. General EqSrc
remains undischarged, and the Distance-to-GR and metric-use ledgers are
unchanged.

The chain-homology route remains \texttt{not\_frozen} for this one packet
because v3 adds a typed cross-complex factorization theorem, the exact
induced-\(H_1\) reflection criterion, epistemically precise fail-closed
branches, and a falsifiable noninvertible witness. Repeating v2, or repeating
v3 without fresh audit or new mathematical payload, would be orbital and
would trigger freeze review. The earlier exact quotient-table cycle remains
locally frozen.

The next lawful packet is one fresh bounded
\texttt{smuggling-auditor@0.2.0} audit of exact v3. It must test typing,
factorization, induced-\(H_1\) injectivity, the noninvertible witness, preserved
v2 scope, primitive selection, target or process import, and physical
overreads. It may not repair, adopt, stress-test, or promote the candidate.

\section{Forbidden Conclusions}

This artifact does not authorize canonical ontology modification, source-law
adoption, physical-admissibility or physical-covariance establishment,
general EqSrc discharge, RetainH or GenH adoption, \(M_{\mathrm{src}}\)
adoption, \(g_{\mathrm{eff}}\) scope expansion, matter coupling, Einstein
equations, robustness beyond declared hypotheses, Distance-to-GR or
metric-use ledger change, benchmark promotion, Gate Chair verdict, completed
derivation, future source-extension impossibility, program-level no-go, or
global theory rejection.

\section{Source Materials}

\begin{itemize}[leftmargin=*]
  \item Handoff 0755 and the RT-20260718-030 exact smuggling audit.
  \item Proposal-only v2 and its surviving theorem core.
  \item The registered Ontology Formalizer role contract v0.2.0.
  \item The source-extension classification checklist and no-target-import
  guard map.
  \item Standard quotient, homology, and finite-dimensional linear-algebra
  arguments reproduced explicitly above.
\end{itemize}

\end{document}
